\scanpage{192}
% archive sheet number: 191

($k$ un anneau) est-il vrai que si $A$ est une hyperalgèbre (\unsure{non commutative?}) sur $k$, alors $A$ \deleted{\unsure{commut.}} est limite inductive de \deleted{\unsure{sous-algèbres de}}\deleted{\unsure{hyperalgèbres}}\inserted{hyperalgèbres} \uline{de type fini} \deleted{\illegible{3}}\inserted{\unsure{du pt de vue associatif}}. Il suffit pour ceci que $A$ soit limite inductive de ss-$k$-modules de type fini $V$ tels que $\Delta V\subset\operatorname{Im}V\otimes V$, et de plus [$A_V$ dés la ss-$k$-algèbre associative engendrée] \deleted{\unsure{et}} $A_V\otimes A_V\mathbin{-}A\otimes A$ soit \uline{injective}. \deleted{\unsure{Ce qui}}
\margin{\diagram{\begin{tikzcd}[ampersand replacement=\&,row sep=small]
G\arrow[r,"\text{épim}"]\&G'\\
A\&A'\arrow[l]
\end{tikzcd}}}

On voit facilement qu'il en est ainsi si $k$ est un \uline{corps} [alors la \deleted{première}\inserted{deuxième} \unsure{cond\textsuperscript{on}} est remplie automatiquement; pour la première, il est possible qu'elle soit vraie \deleted{\unsure{au moins}}\inserted{sous des} \deleted{si $k$ est noethérien:} conditions beaucoup plus générales, telles que: $A$ libre].
\margin{$\unsure{\ast}$}

\uline{Sur un corps} $k$, \deleted{\unsure{un groupe affine}}
\[
\underline{\text{pro (groupe alg. affine)}}=\underline{\text{grp. \unsure{alg.|sch.} affine}}
\]
\deleted{\unsure{$G$}}

Le ``dual'' d'un \deleted{\unsure{tel}} groupe \inserted{comm. sur $k$} satisfaisant la condition (1) est un groupe (ind.fini) $\deleted{\unsure{D}}\widehat G$, sur $S$, dont l'hyperalgèbre est l'algèbre affine de $G$. \unsure{Se bornant au} type fini sur un corps alg. clos, on trouve encore
% slanted
\margin{\unsure{groupe formel} au sens de Cartier}
